% Proposed replacement for the broad higher-branch monotonicity direction
% in 10-discovery.tex. Import the matching labels/citations before use.

Higher zero branches do not have one uniform direction throughout the
Lerch deformation. The earlier boundary, bifurcation and global-phase
results already exhibit turning or increasing branches at small indices.
The diagonal theorem now shows that, for $k=n-1$, the smallest zero
issuing from $a=1$ initially moves right for infinitely many $n$ and
left for infinitely many $n$. This holds for every positive lattice
spacing. Its exact velocity generating function is
\[
 \sum_{n\ge1}v_n(A)t^n=\log A-u_A(t),\qquad
 e^{u_A(t)}=A+t u_A(t),\quad u_A(0)=\log A,\quad A>1.
\]
The fixed-Laurent-index eventual-decrease theorem remains compatible
with this result. If $K_n$ guarantees decrease of every branch for
all $k\ge K_n$, then $K_n>n-1$ for infinitely many $n$.
The resulting questions concern effective index-dependent thresholds,
joint-index transition regimes, and turning-point counts below those
thresholds.

The optimal constant exterior operator is now explicit. If $r_n$ is
the unique positive root of $r^3+nr^2-nr-n=0$, its coefficient is
$\lambda_n=(2-r_n)e^{-n-r_n}$. The least uniform cutoff for positive
geometric weights is attained and yields a strict summed inequality.
For the strict pointwise problem, the same cutoff is an unattained
infimum because of an outer tangency. This resolves the corresponding
incoming research question with its endpoint qualification.

The Gaussian harmonic sums also have an entire order interpolation
$S(s)$. Its real zeros are exactly one simple zero in each
$(-2m-1,-2m)$ for $m\ge1$, and their displacements and spacings have
explicit asymptotic expansions. This analytic result does not resolve
the separate weight-seven $S_6$ period-reduction conjecture.
